% TODO(team): Rehearse the architecture explanation and optional real-media demo.
\section{System}
\begin{frame}{One upload budget serves two demands}
\centering
% Use explicit edge anchors and two-line labels in the gaps between boxes.
\begin{tikzpicture}
\node[box,text width=26mm] (p) at (0,0) {Surviving Peer\\Shared upload cap};
\node[box,text width=27mm] (v) at (6,1.25) {Playback client};
\node[box,text width=27mm] (c) at (6,-1.25) {Controller\\Copy and verify};
\node[box,text width=25mm] (t) at (10,-1.25) {Replacement Peer};
\draw[flow] (p.east)--node[midway,above=2mm,align=center,font=\scriptsize]{foreground\\reads}(v.west);
\draw[flow,draw=Orange] (p.east)--node[midway,below=2mm,align=center,font=\scriptsize]{repair\\reads}(c.west);
\draw[flow,draw=Orange] (c.east)--(t.west);
\end{tikzpicture}
\takeaway{Repair restores a second source, but consumes bandwidth needed by viewers.}
\note{Conceptual data paths, not a measured traffic trace. The surviving source uploads repair bytes to the Controller, which forwards them to a replacement Peer.}
\end{frame}

\begin{frame}{Control plane and data plane}
\centering
\begin{tikzpicture}
\node[box,text width=30mm] (c1) at (0,0) {Controller 1\\Leader};
\node[box,text width=30mm] (c2) at (4.5,0) {Controller 2};
\node[box,text width=30mm] (c3) at (9,0) {Controller 3};
\draw[flow,<->] (c1)--node[midway,above=2mm,font=\scriptsize]{Raft}(c2);
\draw[flow,<->] (c2)--(c3);
\node[box,text width=28mm] (p1) at (0,-2) {Peer 1\\Index + blobs};
\node[box,text width=28mm] (p2) at (4.5,-2) {Peer 2\\Index + blobs};
\node[box,text width=28mm] (p3) at (9,-2) {Peer 3\\Index + blobs};
\draw[meta] (c1)--node[midway,right=2mm,font=\scriptsize]{metadata}(p1);
\draw[meta] (c2)--(p2);\draw[meta] (c3)--(p3);
\end{tikzpicture}
\small
\begin{itemize}
\item Raft: manifests, committed receipts, publication state, persistent tasks.
\item Browser: Peer $\rightarrow$ Gateway $\rightarrow$ Browser.
\item Experiment: Peer $\rightarrow$ headless client; Gateway is outside Raft.
\end{itemize}
\note{Dashed links indicate control information, not video replication through Raft. Any current leader coordinates maintenance; the vertical links do not assign a Peer to a particular Controller. Peers use independent stores and upload limiters.}
\end{frame}

\begin{frame}{Publish only after two durable copies}
\begin{enumerate}
\item Create a \textbf{DRAFT} manifest.
\item Upload segments; Peers validate size and SHA-256 and persist bytes.
\item Confirm receipts from \textbf{two distinct holders per segment}.
\item Commit \textbf{READY}; expose the video for playback.
\end{enumerate}
\takeaway{Raft commits metadata. Peers store the video bytes outside the Raft log.}
\note{The application checks Peer receipts before publication. Browser playback uses continuous fMP4. Experimental payloads are deterministic bytes with simulated media durations.}
\end{frame}

\begin{frame}{Playback recovery $\neq$ redundancy recovery}
\begin{block}{Foreground recovery}
Peer fails $\rightarrow$ try another holder $\rightarrow$ verify bytes $\rightarrow$ continue delivery.
\end{block}
\begin{block}{Background repair}
Detect shortage $\rightarrow$ grace interval $\rightarrow$ commit task\\
$\rightarrow$ copy and verify $\rightarrow$ commit replacement registration.
\end{block}
\takeaway{Continuing playback does not prove that redundancy has been restored.}
\note{A repaired segment is available from the replacement only after successful installation and registration. If both original holders stop before repair reaches a needed segment, that segment remains unavailable.}
\end{frame}

\begin{frame}{Application implementation and validation}
\small
\begin{tabularx}{\textwidth}{YY}
\toprule Application work & Foundation / tool \\\midrule
Publication and persistent task lifecycle & HashiCorp Raft, BoltDB, snapshots\\
Durable indexes and immutable objects & Go HTTP / filesystem, SHA-256\\
Playback, budgets and experiment harness & FFmpeg / FFprobe for media gates\\
Console: faults, cleanup and video deletion & React and TypeScript\\\bottomrule
\end{tabularx}
\vspace{3mm}
Archived gates: build, static/race, smoke, repair, cleanup, deletion, exact fragment timing and media decoding.
\evidence{Evidence: results/compact-study/validation/gates.json.}
\note{The archived media gates do not establish browser performance in the headless study. Attribute the consensus implementation to HashiCorp rather than presenting it as our own Raft implementation.}
\end{frame}

\ifshowdemo
\begin{frame}{Optional demonstration: from upload to playback}
\begin{enumerate}
\item Start three Controllers and three Peers.
\item Upload a rehearsed short clip; observe packaging and publication.
\item Play the video; inspect per-Peer segment counts.
\item Return to the controlled experimental results.
\end{enumerate}
\takeaway{A functional demonstration and comparative performance evidence are different.}
\note{Prepare an already imported clip and a recording as fallbacks. Avoid an unplanned live failure. No screenshot or recording is fabricated or embedded in this deck.}
\end{frame}
\fi
